%% Question bank chapter file - generated by bank_tools.py build (do not edit while the bank is still being built:
%% it is regenerated from the per-test data in Question Bank\data). Every question is one bankq block:
%%   \begin{bankq}[id=..., source={...}, paper=.., qno=.., type=..., subtopic={...}, answer={...}, answer_src=key|solved,
%%                 working={...}, difficulty=1-10, diff_parts={K S I C F}, diff_note={...},
%%                 group=<passage id or empty>, also_in={duplicate copies in other tests}]
%% difficulty = min(10, K+S+I+C+F-1): K concepts 1-3, S steps 1-4, I insight 0-2, C computation 0-2, F format 0-1
%%   ...question LaTeX, options, figures (\bfig / \bfigi -> img/)...
%%   \end{bankq}
%% Passages shared by several questions are bankpassage blocks; their questions carry group=<passage id>.
\documentclass[10pt]{article}
\usepackage[a4paper,margin=1.8cm]{geometry}
\usepackage[T1]{fontenc}
\usepackage{amsmath,amssymb,mathtools}
\usepackage{graphicx}
\usepackage[export]{adjustbox}
\usepackage{xcolor,needspace,keyval}
\usepackage{enumitem}
\usepackage{booktabs,array,tabularx}
\usepackage[hidelinks]{hyperref}
\setlength{\parindent}{0pt}
\setlength{\parskip}{3pt}
\makeatletter
\define@key{bq}{id}{\def\bq@id{#1}}\define@key{bq}{source}{\def\bq@src{#1}}\define@key{bq}{paper}{\def\bq@paper{#1}}
\define@key{bq}{qno}{\def\bq@qno{#1}}\define@key{bq}{type}{\def\bq@type{#1}}\define@key{bq}{subtopic}{\def\bq@sub{#1}}
\define@key{bq}{answer}{\def\bq@ans{#1}}\define@key{bq}{answer_src}{\def\bq@asrc{#1}}\define@key{bq}{working}{\def\bq@work{#1}}
\define@key{bq}{group}{\def\bq@group{#1}}\define@key{bq}{also_in}{\def\bq@also{#1}}
\define@key{bq}{difficulty}{\def\bq@diff{#1}}\define@key{bq}{diff_parts}{\def\bq@dparts{#1}}\define@key{bq}{diff_note}{\def\bq@dnote{#1}}
\newenvironment{bankq}[1][]{\def\bq@work{}\def\bq@also{}\def\bq@group{}\def\bq@diff{?}\def\bq@dparts{}\def\bq@dnote{}\setkeys{bq}{#1}%
  \par\needspace{4\baselineskip}\medskip\noindent{\small\textbf{\bq@id}\hfill\color{gray}[\bq@src, Paper \bq@paper, Q\bq@qno]
  \bq@type\ $\cdot$ \bq@sub\ $\cdot$ \textbf{D\bq@diff}}\par\nopagebreak\relax}
  {\par{\small\color{teal}Answer: \textbf{\bq@ans}\ifx\bq@work\@empty\else\ (solved: \bq@work)\fi
  \ifx\bq@also\@empty\else\par\color{gray}Also in: \bq@also\fi}\par}
\define@key{bp}{id}{\def\bp@id{#1}}\define@key{bp}{source}{\def\bp@src{#1}}\define@key{bp}{questions}{\def\bp@qs{#1}}
\newenvironment{bankpassage}[1][]{\setkeys{bp}{#1}\par\needspace{4\baselineskip}\bigskip\noindent
  {\small\textbf{Passage \bp@id}\hfill\color{gray}[\bp@src; questions \bp@qs]}\par\nopagebreak\relax\itshape}{\par}
\makeatother
\newcommand{\bfig}[2]{\par\smallskip{\centering\includegraphics[width=#1pt,max width=0.62\linewidth,max height=0.32\textheight]{img/#2}\par}\smallskip}
\newcommand{\bfigi}[2]{\includegraphics[width=#1pt,max width=\linewidth,max height=0.32\textheight,valign=c]{img/#2}}
\newcommand{\bmissing}[1]{\textit{[figure missing in the source conversion: \detokenize{#1}]}}
\newenvironment{qoptions}{\begin{itemize}[leftmargin=1.8em,itemsep=1pt,topsep=2pt,label={}]}{\end{itemize}}
\newcommand{\opt}[2]{\item[(#1)] #2}
\providecommand{\fwstart}{}\providecommand{\fwend}{}\providecommand{\qfullsplit}{}
\newcolumntype{L}[1]{>{\raggedright\arraybackslash}p{#1}}
\newcolumntype{C}[1]{>{\centering\arraybackslash}p{#1}}
\newcolumntype{R}[1]{>{\raggedleft\arraybackslash}p{#1}}
\providecommand{\sectionhdr}[1]{}\providecommand{\subsectionhdr}[1]{}


\begin{document}
{\Large\bfseries d \& f block --- Question Bank}\par\smallskip
{\small Chemistry chapter C37. Questions from all converted test papers, grouped by source series. Answers: printed key where it matches the paper, otherwise solved (answer\_src=solved, with working). Exact repeats across papers appear once, with also\_in listing the other copies.}\par\medskip\hrule

\section*{Score-I papers}
\begin{bankq}[id=SI12-P1-CHEM-S1-Q10, source={Score-I - Test 12 (09 Mar 2022)}, paper=1, qno=10, type=multi, subtopic={Transition elements}, answer={A,C,D}, answer_src=key, difficulty=2, diff_parts={K1 S1 I0 C0 F1}, diff_note={concepts: CuSO4 and excess KI | steps: Cu2+ is reduced to white CuI and I2 forms I3-; check the four statements}]
Consider the following transformation:\\
CuSO$_{4(\text{aq.})}$ + KI$_{(\text{excess})}$ $\to$ Product\\
Select the INCORRECT statement.
\begin{qoptions}
\opt{A}{Product contains [Cu(H$_2$O)$_4$]$^{2+}$ ions.}
\opt{B}{Presence of brown colour in product is due to I$_3^-$ ion.}
\opt{C}{Oxidation state of sulphur in reactant and product is different}
\opt{D}{White ppt. of CuI$_2$ is observed in product}
\end{qoptions}
\end{bankq}

\section*{TOAS papers}
\begin{bankq}[id=TOAS17-P2-CHEM-S2-Q09, source={TOAS - Test 17 (18 Sep 2022)}, paper=2, qno=9, type=single, subtopic={Transition elements}, answer={213}, answer_src=solved, working={$\mathrm{Mn_2O_7}$ is a green oil (2), ammonium dichromate is orange (1), chromyl chloride is reddish brown (3)}, difficulty=2, diff_parts={K1 S2 I0 C0 F0}, diff_note={concepts: colours of compounds | steps: Mn2O7 green; ammonium dichromate orange; chromyl chloride reddish brown}]
Match the colour with given compounds
\begin{center}\footnotesize\setlength{\tabcolsep}{2pt}
\begin{tabular}{|L{0.35cm}L{1.8cm}|L{0.35cm}L{1.8cm}|}\hline
\multicolumn{2}{|c|}{List-I (Compounds)}&\multicolumn{2}{c|}{List-II (Colour)}\\\hline
P&Mn$_2$O$_7$&1&Orange\\\hline
Q&(NH$_4$)$_2$Cr$_2$O$_7$&2&Green\\\hline
R&CrO$_2$Cl$_2$&3&Reddish Brown\\\hline
\end{tabular}\end{center}
Match list-I and list-II and give your answer as abc where a,b and c are the match of 1, 2, 3 from list-II.

\textbf{(for example Matching of P is = 2, Q = 3 \& R = 1 then your answer is 231)}
\end{bankq}
\begin{bankq}[id=TOAS18-P1-CHEM-S1-Q01, source={TOAS - Test 18 (30 Oct 2022)}, paper=1, qno=1, type=multi, subtopic={Transition elements}, answer={A,B,C,D}, answer_src=solved, working={Potassium manganate is green, potassium superoxide is yellow, chromate is yellow and lead dioxide is brown: all coloured}, difficulty=3, diff_parts={K1 S2 I0 C0 F1}, diff_note={concepts: colours of compounds | steps: manganate green; superoxide yellow; chromate yellow; lead dioxide brown}]
Which of the following is/are coloured?
\begin{qoptions}
\opt{A}{K$_2$MnO$_4$}
\opt{B}{KO$_2$}
\opt{C}{K$_2$CrO$_4$}
\opt{D}{PbO$_2$}
\end{qoptions}
\end{bankq}
\end{document}
